\uuid{YcO5}
\titre{Matrice d'une application linéaire/Noyau/Image/Rang}
\auteur{hanani ; mzari}
\organisation{exo7}
\licence{CC BY-NC-SA 4.0}
\datecreate{2020-01-20}
\niveau{L1}
\module{Algèbre}
\chapitre{Matrice}
\sousChapitre{Matrice et application linéaire}
\difficulte{3}
\theme{Matrice d'une application linéaire, Noyau, Image, Rang}
\exo7qcmid{lille-2.matrice-application.025}

\contenu{
\begin{question}
\qclassification{108.03}
\qtags{motcle=Matrice d'une application linéaire/Noyau/Image/Rang}
Soit $E$ un espace vectoriel de dimension $3$ et $f$ un endomorphisme non nul de $E$ tel que $f^2=0$. Quelles sont les assertions vraies ?
\begin{answers}
\good{$\Im f \subset \ker f$}.
\bad{$\Im f = \ker f$}.
\bad{Le rang de $f$ est $2$}.
\good{Il existe une base de $E$ dans laquelle le matrice de $f$ est : $\left(\begin{array}{rcc}0&0&a\\
0&0&0\\ 0&0&0\\\end{array}\right)$, où $a$ est un réel non nul}.
\end{answers}
\begin{explanations} Comme $f^2=0$, $\Im f \subset \ker f$.
D'après le théorème du rang, on déduit que $\dim \ker f=2$ et $\mbox{rg} (f)=1$.
\vskip0mm
Soit $\{u\}$ une base de $\Im f$. On complète cette base pour obtenir une base $\{u,v\}$ de $\ker f$, puis, on complète cette
dernière base pour obtenir une base $\{u,v,w\}$ de $E$. Alors, la matrice de $f$ dans cette base est de la forme :
$\left(\begin{array}{rcc}
0&0&a\\
0&0&0\\
0&0&0\\
\end{array}\right)$, où $a$ est un réel non nul.
\end{explanations}
\end{question}
}
